\RequirePackage[T1]{fontenc}
\documentclass[10pt,conference]{IEEEtran}
\usepackage{amsmath,amssymb,amsthm,graphicx,booktabs,array,url,flushend}
\usepackage[hidelinks]{hyperref}
\usepackage{tikz}
\usetikzlibrary{arrows.meta,positioning}
\emergencystretch=2em
\newtheorem{theorem}{Theorem}
\newtheorem{lemma}{Lemma}
\newtheorem{proposition}{Proposition}
\theoremstyle{definition}
\newtheorem{definition}{Definition}
\newtheorem{remarknum}{Remark}
\newcommand{\code}[1]{\texttt{#1}}
\newcommand{\G}[1]{G#1}
\newcommand{\tool}{ReleaseWeave}
\newcommand{\lang}{RIL}
\begin{document}
\title{Recovering Versioned Releases Across Build and Publish Boundaries}
\maketitle
\begin{abstract}
Distributed releases must transfer dependency outputs and preserve unfinished
work after provider success. Our protocol combines deterministic planning,
stage-aware scheduling, and consumer-owned release obligations.
Release planning is a pure function of fixed history, records, graph, and
configuration. An accompanying agent guide lets coding agents inspect plans
and invoke CI releases through this deterministic engine.
A proposed execution profile adds workers to repositories and polyrepos.
Nodes share a schema and default to orchestrator capability; explicit workers
cannot initiate releases. The CI-entry node locks participating repositories
and schedules tasks within concurrency limits. Temporary Git branches and
verified outputs enable remote reuse without rebuilding providers; publication retains dependency gates
and source-owned records.
The single-history model gives eight conditional mathematical recovery guarantees.
All 120 controlled fault episodes recover after repair; an analytical
source-discharge alternative loses unfinished work in 84 (70\%). This compares
admission rules, not competing tools or production failure rates.
For LLVM on 6--10 compatible 16-core machines, an illustrative neutral
scenario projects about 103 minutes against a 154-minute upstream Linux
release baseline (33\% less),
conditional on a wider task graph, sharded tests, and five-minute output hops.
These are forecasts, not measured distributed results; the existing
four-package graph offers little parallelism. Full implementation refinement
and external exactly-once effects remain outside the claims.
\end{abstract}
\begin{IEEEkeywords}
software release engineering, continuous integration and delivery,
dependency scheduling, recovery, idempotence
\end{IEEEkeywords}
\input{introduction}
\input{system}
\input{model}
\input{guarantees}
\input{evaluation}
\input{discussion}
\clearpage
\IEEEtriggeratref{34}
\input{references}
\end{document}
